\section{AHSME 1966}
        \begin{enumerate}
        	\item (\textit{AHSME 1966}) Given that the ratio of $3x - 4$ to $y + 15$ is constant, and $y = 3$ when $x = 2$, then, when $y = 12$, $x$ equals:


 $\text{(A)} \ \frac 18 \qquad \text{(B)} \ \frac 73 \qquad \text{(C)} \ \frac78 \qquad \text{(D)} \ \frac72 \qquad \text{(E)} \ 8$


 

	\item (\textit{AHSME 1966}) When the base of a triangle is increased 10\% and the altitude to this base is decreased 10\%, the change in area is 


 $\text{(A)}\ 1\%~\text{increase}\qquad\text{(B)}\ \frac{1}2\%~\text{increase}\qquad\text{(C)}\ 0\%\qquad\text{(D)}\ \frac{1}2\% ~\text{decrease}\qquad\text{(E)}\ 1\% ~\text{decrease}$


 

	\item (\textit{AHSME 1966}) If the arithmetic mean of two numbers is $6$ and their geometric mean is $10$, then an equation with the given two numbers as roots is:


 $\text{(A)} \ x^2 + 12x + 100 = 0 ~~ \text{(B)} \ x^2 + 6x + 100 = 0 ~~ \text{(C)} \ x^2 - 12x - 10 = 0$
$\text{(D)} \ x^2 - 12x + 100 = 0 \qquad \text{(E)} \ x^2 - 6x + 100 = 0$


 

	\item (\textit{AHSME 1966}) Circle I is circumscribed about a given square and circle II is inscribed in the given square. If $r$ is the ratio of the area of circle $I$ to that of circle $II$, then $r$ equals:


 $\text{(A)}  \sqrt 2 \qquad \text{(B)}  2 \qquad \text{(C)}  \sqrt 3 \qquad \text{(D)}  2\sqrt 2 \qquad \text{(E)}  2\sqrt 3$


 

	\item (\textit{AHSME 1966}) The number of values of $x$ satisfying the equation


 $\frac {2x^2 - 10x}{x^2 - 5x} = x - 3$


 is:


 $\text{(A)} \ \text{zero} \qquad \text{(B)} \ \text{one} \qquad \text{(C)} \ \text{two} \qquad \text{(D)} \ \text{three} \qquad \text{(E)} \ \text{an integer greater than 3}$


 

	\item (\textit{AHSME 1966}) $AB$ is the diameter of a circle centered at $O$. $C$ is a point on the circle such that angle $BOC$ is $60^\circ$. If the diameter of the circle is $5$ inches, the length of chord $AC$, expressed in inches, is:


 $\text{(A)} \ 3 \qquad \text{(B)} \ \frac {5\sqrt {2}}{2} \qquad \text{(C)} \frac {5\sqrt3}{2} \ \qquad \text{(D)} \ 3\sqrt3 \qquad \text{(E)} \ \text{none of these}$


 

	\item (\textit{AHSME 1966}) Let $\frac {35x - 29}{x^2 - 3x + 2} = \frac {N_1}{x - 1} + \frac {N_2}{x - 2}$ be an identity in $x$. The numerical value of $N_1N_2$ is:


 $\text{(A)} \ - 246 \qquad \text{(B)} \ - 210 \qquad \text{(C)} \ - 29 \qquad \text{(D)} \ 210 \qquad \text{(E)} \ 246$


 

	\item (\textit{AHSME 1966}) The length of the common chord of two intersecting circles is $16$ feet. If the radii are $10$ feet and $17$ feet, a possible value for the distance between the centers of the circles, expressed in feet, is:


 $\text{(A)} \ 27 \qquad \text{(B)} \ 21 \qquad \text{(C)} \ \sqrt {389} \qquad \text{(D)} \ 15 \qquad \text{(E)} \ \text{undetermined}$


 

	\item (\textit{AHSME 1966}) If $x = (\log_82)^{(\log_28)}$, then $\log_3x$ equals:


 $\text{(A)} \ - 3 \qquad \text{(B)} \ - \frac13 \qquad \text{(C)} \ \frac13 \qquad \text{(D)} \ 3 \qquad \text{(E)} \ 9$


 

	\item (\textit{AHSME 1966}) If the sum of two numbers is 1 and their product is 1, then the sum of their cubes is:


 $\text{(A)} \ 2 \qquad \text{(B)} \ - 2 - \frac {3i\sqrt {3}}{4} \qquad \text{(C)} \ 0 \qquad \text{(D)} \ - \frac {3i\sqrt {3}}{4} \qquad \text{(E)} \ - 2$


 

	\item (\textit{AHSME 1966}) The sides of triangle $BAC$ are in the ratio $2: 3: 4$. $BD$ is the angle-bisector drawn to the shortest side $AC$, dividing it into segments $AD$ and $CD$. If the length of $AC$ is $10$, then the length of the longer segment of $AC$ is:


 $\text{(A)} \ 3\frac12 \qquad \text{(B)} \ 5 \qquad \text{(C)} \ 5\frac57 \qquad \text{(D)} \ 6 \qquad \text{(E)} \ 7\frac12$


 

	\item (\textit{AHSME 1966}) The number of real values of $x$ that satisfy the equation 
 \[(2^{6x+3})(4^{3x+6})=8^{4x+5}\] is:


 $\text{(A)  zero} \qquad \text{(B)  one} \qquad \text{(C)  two} \qquad \text{(D)  three} \qquad \text{(E)  greater than 3}$


 

	\item (\textit{AHSME 1966}) The number of points with positive rational coordinates selected from the set of points in the $xy$-plane such that $x+y \le 5$, is:


 $\text{(A)} \ 9 \qquad \text{(B)} \ 10 \qquad \text{(C)} \ 14 \qquad \text{(D)} \ 15 \qquad \text{(E) infinite}$


 

	\item (\textit{AHSME 1966}) The length of rectangle $ABCD$ is 5 inches and its width is 3 inches. Diagonal $AC$ is divided into three equal segments by points $E$ and $F$. The area of triangle $BEF$, expressed in square inches, is:


 $\text{(A)} \frac{3}{2} \qquad \text{(B)} \frac {5}{3} \qquad \text{(C)} \frac{5}{2} \qquad \text{(D)} \frac{1}{3}\sqrt{34} \qquad \text{(E)} \frac{1}{3}\sqrt{68}$


 

	\item (\textit{AHSME 1966}) If $x-y>x$ and $x+y<y$, then


 $\text{(A) } y<x \quad \text{(B) } x<y \quad \text{(C) } x<y<0 \quad \text{(D) } x<0,y<0 \quad \text{(E) } x<0,y>0$


 

	\item (\textit{AHSME 1966}) If $\frac{4^x}{2^{x+y}}=8$ and $\frac{9^{x+y}}{3^{5y}}=243$, $x$ and $y$ real numbers, then $xy$ equals:


 $\text{(A) } \frac{12}{5} \quad \text{(B) } 4 \quad \text{(C) } 6 \quad \text{(D)} 12 \quad \text{(E) } -4$


 

	\item (\textit{AHSME 1966}) The number of distinct points common to the curves $x^2+4y^2=1$ and $4x^2+y^2=4$ is:


 $\text{(A) } 0 \quad \text{(B) } 1 \quad \text{(C) } 2 \quad \text{(D) } 3 \quad \text{(E) } 4$


 

	\item (\textit{AHSME 1966}) In a given arithmetic sequence the first term is $2$, the last term is $29$, and the sum of all the terms is $155$. The common difference is:


 $\text{(A) } 3 \qquad \text{(B) } 2 \qquad \text{(C) } \frac{27}{19} \qquad \text{(D) } \frac{13}{9} \qquad \text{(E) } \frac{23}{38}$


 

	\item (\textit{AHSME 1966}) Let $s_1$ be the sum of the first $n$ terms of the arithmetic sequence $8,12,\cdots$ and let $s_2$ be the sum of the first $n$ terms of the arithmetic sequence $17,19,\cdots$. Assume $n \ne 0$. Then $s_1=s_2$ for:


 $\text{(A) no value of } n \quad \text{(B) one value of } n \quad \text{(C) two values of } n \quad \text{(D) four values of } n \quad \text{(E) more than four values of } n$


 

	\item (\textit{AHSME 1966}) The negation of the proposition "For all pairs of real numbers $a,b$, if $a=0$, then $ab=0$" is: There are real numbers $a,b$ such that


 $\text{(A) } a\ne 0 \text{ and } ab\ne 0 \qquad \text{(B) } a\ne 0 \text{ and } ab=0  \qquad \text{(C) } a=0 \text{ and } ab\ne 0$


 $\text{(D) } ab\ne 0 \text{ and } a\ne 0 \qquad \text{(E) } ab=0 \text{ and } a\ne 0$


 

	\item (\textit{AHSME 1966}) 
                    \begin{figure}[H]
                        \centering
                        \includegraphics[width=0.5\textwidth]{Problem/AHSME_I_1966_21_Problem_Figure_1.png}
                    \end{figure}
                    


 An "$n$-pointed star" is formed as follows: the sides of a convex polygon are numbered consecutively $1,2,\cdots ,k,\cdots,n,\text{ }n\ge 5$; for all $n$ values of $k$, sides $k$ and $k+2$ are non-parallel, sides $n+1$ and $n+2$ being respectively identical with sides $1$ and $2$; prolong the $n$ pairs of sides numbered $k$ and $k+2$ until they meet. (A figure is shown for the case $n=5$).


 Let $S$ be the degree-sum of the interior angles at the $n$ points of the star; then $S$ equals:


 $\text{(A) } 180 \quad \text{(B) } 360 \quad \text{(C) } 180(n+2) \quad \text{(D) } 180(n-2) \quad \text{(E) } 180(n-4)$


 

	\item (\textit{AHSME 1966}) Consider the statements: (I)$\sqrt{a^2+b^2}=0$,   (II) $\sqrt{a^2+b^2}=ab$,   (III) $\sqrt{a^2+b^2}=a+b$,   (IV) $\sqrt{a^2+b^2}=a - b$, where we allow $a$ and $b$ to be real or complex numbers. Those statements for which there exist solutions other than $a=0$ and $b=0$, are:


 $\text{(A)  (I),(II),(III),(IV)}\quad \text{(B)  (II),(III),(IV) only} \quad \text{(C)  (I),(III),(IV) only}\quad \text{(D)  (III),(IV) only} \quad \text{(E)  (I) only}$


 

	\item (\textit{AHSME 1966}) If $x$ is real and $4y^2+4xy+x+6=0$, then the complete set of values of $x$ for which $y$ is real, is:


 $\text{(A) } x\le-2 \text{ or } x\ge3 \quad \text{(B) }  x\le2 \text{ or } x\ge3 \quad \text{(C) }  x\le-3 \text{ or } x\ge2 \quad \\ \text{(D) } -3\le x\le2 \quad \text{(E) } -2\le x\le3$


 

	\item (\textit{AHSME 1966}) If $Log_M{N}=Log_N{M},M \ne N,MN>0,M \ne 1, N \ne 1$, then $MN$ equals:


 $\text{(A) } \frac{1}{2} \quad \text{(B) } 1 \quad \text{(C) } 2 \quad \text{(D) } 10 \\ \text{(E) a number greater than 2 and less than 10}$


 

	\item (\textit{AHSME 1966}) If $F(n+1)=\frac{2F(n)+1}{2}$ for $n=1,2,\cdots$ and $F(1)=2$, then $F(101)$ equals:


 $\text{(A) } 49 \quad \text{(B) } 50 \quad \text{(C) } 51 \quad \text{(D) } 52 \quad \text{(E) } 53$


 

	\item (\textit{AHSME 1966}) Let $m$ be a positive integer and let the lines $13x+11y=700$ and $y=mx-1$ intersect in a point whose coordinates are integers. Then m can be:


 $\text{(A) 4 only} \quad \text{(B) 5 only} \quad \text{(C) 6 only} \quad \text{(D) 7 only} \\ \text{(E) one of the integers 4,5,6,7 and one other positive integer}$


 

	\item (\textit{AHSME 1966}) At his usual rate a man rows 15 miles downstream in five hours less time than it takes him to return. If he doubles his usual rate, the time downstream is only one hour less than the time upstream. In miles per hour, the rate of the stream's current is:


 $\text{(A) } 2 \quad \text{(B) } \frac{5}{2} \quad \text{(C) } 3 \quad \text{(D) } \frac{7}{2} \quad \text{(E) } 4$


 

	\item (\textit{AHSME 1966}) Five points $O,A,B,C,D$ are taken in order on a straight line with distances $OA = a$, $OB = b$, $OC = c$, and $OD = d$.  $P$ is a point on the line between $B$ and $C$ and such that $AP: PD = BP: PC$.  Then $OP$ equals:


 $\textbf{(A)} \frac {b^2 - bc}{a - b + c - d} \qquad \textbf{(B)} \frac {ac - bd}{a - b + c - d} \\  \textbf{(C)} - \frac {bd + ac}{a - b + c - d} \qquad \textbf{(D)} \frac {bc + ad}{a + b + c + d} \qquad \textbf{(E)} \frac {ac - bd}{a + b + c + d}$


 

	\item (\textit{AHSME 1966}) The number of positive integers less than $1000$ divisible by neither $5$ nor $7$ is:


 $\text{(A) } 688 \quad \text{(B) } 686 \quad \text{(C) } 684 \quad \text{(D) } 658 \quad \text{(E) } 630$


 

	\item (\textit{AHSME 1966}) If three of the roots of $x^4+ax^2+bx+c=0$ are $1$, $2$, and $3$ then the value of $a+c$ is:


 $\text{(A) } 35 \quad \text{(B) } 24 \quad \text{(C) } -12 \quad \text{(D) } -61 \quad \text{(E) } -63$


 

	\item (\textit{AHSME 1966}) 
                    \begin{figure}[H]
                        \centering
                        \includegraphics[width=0.5\textwidth]{Problem/AHSME_I_1966_31_Problem_Figure_1.png}
                    \end{figure}
                    
Triangle $ABC$ is inscribed in a circle with center $O'$. A circle with center $O$ is inscribed in triangle $ABC$. $AO$ is drawn, and extended to intersect the larger circle in $D$. Then we must have:


 $\text{(A) } CD=BD=O'D \quad \text{(B) } AO=CO=OD \quad \text{(C) } CD=CO=BD \\ \text{(D) } CD=OD=BD \quad \text{(E) } O'B=O'C=OD$


 

	\item (\textit{AHSME 1966}) Let $M$ be the midpoint of side $AB$ of triangle $ABC$. Let $P$ be a point on $AB$ between $A$ and $M$, and let $MD$ be drawn parallel to $PC$ and intersecting $BC$ at $D$. If the ratio of the area of triangle $BPD$ to that of triangle $ABC$ is denoted by $r$, then


 $\text{(A) } \frac{1}{2}<r<1 \text{,  depending upon the position of P} \\ \text{(B) } r=\frac{1}{2} \text{,  independent of the position of P} \\ \text{(C) } \frac{1}{2} \le r <1 \text{,  depending upon the position of P} \\ \text{(D) } \frac{1}{3}<r<\frac{2}{3} \text{,  depending upon the position of P}\\ \text{(E) } r=\frac{1}{3} \text{,  independent of the position of P}$


 

	\item (\textit{AHSME 1966}) If $ab \ne 0$ and $|a| \ne |b|$, the number of distinct values of $x$ satisfying the equation


 
 \[\frac{x-a}{b}+\frac{x-b}{a}=\frac{b}{x-a}+\frac{a}{x-b},\]


 is:


 $\text{(A) zero}  \quad \text{(B) one}  \quad \text{(C) two}  \quad \text{(D) three}  \quad \text{(E) four}$


 

	\item (\textit{AHSME 1966}) Let $r$ be the speed in miles per hour at which a wheel, $11$ feet in circumference, travels. If the time for a complete rotation of the wheel is shortened by $\frac{1}{4}$ of a second, the speed $r$ is increased by $5$ miles per hour. Then $r$ is:


 $\text{(A) } 9 \quad \text{(B) } 10 \quad \text{(C) } 10\frac{1}{2} \quad \text{(D) } 11 \quad \text{(E) } 12$


 

	\item (\textit{AHSME 1966}) Let $O$ be an interior point of triangle $ABC$, and let $s_1=OA+OB+OC$. If $s_2=AB+BC+CA$, then


 $\text{(A) for every triangle } s_2>2s_1,s_1 \le s_2 \\  \text{(B) for every triangle } s_2>2s_1,s_1 < s_2 \\  \text{(C) for every triangle } s_1> \frac{1}{2}s_2,s_1 < s_2 \\  \text{(D) for every triangle } s_2\ge 2s_1,s_1 \le s_2 \\  \text{(E) neither (A) nor (B) nor (C) nor (D) applies to every triangle}$


 

	\item (\textit{AHSME 1966}) Let $(1+x+x^2)^n=a_0+a_1x+a_2x^2+ \cdots + a_{2n}x^{2n}$ be an identity in $x$. If we let $s=a_0+a_2+a_4+\cdots +a_{2n}$, then $s$ equals:


 $\text{(A) } 2^n \quad \text{(B) } 2^n+1 \quad \text{(C) } \frac{3^n-1}{2} \quad \text{(D) } \frac{3^n}{2} \quad \text{(E) } \frac{3^n+1}{2}$


 

	\item (\textit{AHSME 1966}) Three men, Alpha, Beta, and Gamma, working together, do a job in 6 hours less time than Alpha alone, in 1 hour less time than Beta alone, and in one-half the time needed by Gamma when working alone. Let $h$ be the number of hours needed by Alpha and Beta, working together, to do the job. Then $h$ equals:


 $\text{(A) } \frac{5}{2} \quad \text{(B) } \frac{3}{2} \quad \text{(C) } \frac{4}{3} \quad \text{(D) } \frac{5}{4} \quad \text{(E) } \frac{3}{4}$


 

	\item (\textit{AHSME 1966}) In triangle $ABC$ the medians $AM$ and $CN$ to sides $BC$ and $AB$, respectively, intersect in point $O$. $P$ is the midpoint of side $AC$, and $MP$ intersects $CN$ in $Q$. If the area of triangle $OMQ$ is $n$, then the area of triangle $ABC$ is:


 $\text{(A) } 16n \quad \text{(B) } 18n \quad \text{(C) } 21n \quad \text{(D) } 24n \quad \text{(E) } 27n$


 

	\item (\textit{AHSME 1966}) In base $R_1$ the expanded fraction $F_1$ becomes $.373737\cdots$, and the expanded fraction $F_2$ becomes $.737373\cdots$. In base $R_2$ fraction $F_1$, when expanded, becomes $.252525\cdots$, while the fraction $F_2$ becomes $.525252\cdots$. The sum of $R_1$ and $R_2$, each written in the base ten, is:


 $\text{(A) } 24 \quad \text{(B) } 22 \quad \text{(C) } 21 \quad \text{(D) } 20 \quad \text{(E) } 19$


 

	\item (\textit{AHSME 1966}) 
                    \begin{figure}[H]
                        \centering
                        \includegraphics[width=0.5\textwidth]{Problem/AHSME_I_1966_40_Problem_Figure_1.png}
                    \end{figure}
                    
In this figure $AB$ is a diameter of a circle, centered at $O$, with radius $a$. A chord $AD$ is drawn and extended to meet the tangent to the circle at $B$ in point $C$. Point $E$ is taken on $AC$ so the $AE=DC$. Denoting the distances of $E$ from the tangent through $A$ and from the diameter $AB$ by $x$ and $y$, respectively, we can deduce the relation:


 $\text{(A) } y^2=\frac{x^3}{2a-x} \quad \text{(B) } y^2=\frac{x^3}{2a+x} \quad \text{(C) } y^4=\frac{x^2}{2a-x} \\ \text{(D) } x^2=\frac{y^2}{2a-x} \quad \text{(E) } x^2=\frac{y^2}{2a+x}$


 

\end{enumerate}